\begin{ex}
Tính: $36 + 9$.
\shortans{45}
\loigiai{
Ta có $36 + 9 = 45$.
}
\end{ex}

\begin{ex}
Tính: $48 + 7$.
\shortans{55}
\loigiai{
Ta có $48 + 7 = 55$.
}
\end{ex}

\begin{ex}
Tính: $27 + 18$.
\shortans{45}
\loigiai{
Ta có $27 + 18 = 45$.
}
\end{ex}

\begin{ex}
Tính: $35 + 29$.
\shortans{64}
\loigiai{
Ta có $35 + 29 = 64$.
}
\end{ex}

\begin{ex}
Tính: $58 + 36$.
\shortans{94}
\loigiai{
Ta có $58 + 36 = 94$.
}
\end{ex}

\begin{ex}
Tính: $49 + 43$.
\shortans{92}
\loigiai{
Ta có $49 + 43 = 92$.
}
\end{ex}

\begin{ex}
Tính giá trị biểu thức: $28 + 14 + 18$.
\shortans{60}
\loigiai{
Ta có $28 + 14 = 42$, rồi $42 + 18 = 60$.
}
\end{ex}

\begin{ex}
Tính giá trị biểu thức: $37 + 25 + 18$.
\shortans{80}
\loigiai{
Ta có $37 + 25 = 62$, rồi $62 + 18 = 80$.
}
\end{ex}

\begin{ex}
Tính giá trị biểu thức: $46 + 19 + 25$.
\shortans{90}
\loigiai{
Ta có $46 + 19 = 65$, rồi $65 + 25 = 90$.
}
\end{ex}

\begin{ex}
Điền số thích hợp: $38 + 25 = 40 + \dots$
\shortans{23}
\loigiai{
Vế trái: $38 + 25 = 63$. Số cần điền: $63 - 40 = 23$.
}
\end{ex}

\begin{ex}
Điền số thích hợp: $47 + 16 = 50 + \dots$
\shortans{13}
\loigiai{
Vế trái: $47 + 16 = 63$. Số cần điền: $63 - 50 = 13$.
}
\end{ex}

\begin{ex}
Điền số thích hợp: $29 + \dots = 62$.
\shortans{33}
\loigiai{
Số cần điền: $62 - 29 = 33$.
}
\end{ex}

\begin{ex}
Điền số thích hợp: $56 + \dots = 81$.
\shortans{25}
\loigiai{
Số cần điền: $81 - 56 = 25$.
}
\end{ex}

\begin{ex}
Hà có 27 chiếc kẹp tóc, mẹ mua thêm cho Hà 15 chiếc nữa. Hà có tất cả bao nhiêu chiếc kẹp tóc?
\shortans{42}
\loigiai{
Hà có: $27 + 15 = 42$ chiếc kẹp tóc.
}
\end{ex}

\begin{ex}
Trong thư viện trường có 48 cuốn truyện tranh và 36 cuốn sách khoa học. Có tất cả bao nhiêu cuốn sách và truyện?
\shortans{84}
\loigiai{
Có tất cả: $48 + 36 = 84$ cuốn.
}
\end{ex}

\begin{ex}
Một đội trồng rừng đã trồng được 57 cây thông và 28 cây bạch đàn. Đội đã trồng được bao nhiêu cây?
\shortans{85}
\loigiai{
Đội trồng được: $57 + 28 = 85$ cây.
}
\end{ex}

\begin{ex}
Thùng thứ nhất đựng 38 lít dầu, thùng thứ hai đựng nhiều hơn thùng thứ nhất 17 lít dầu. Thùng thứ hai đựng bao nhiêu lít dầu?
\shortans{55}
\loigiai{
Thùng thứ hai đựng: $38 + 17 = 55$ lít.
}
\end{ex}

\begin{ex}
Bao gạo cân nặng 46 kg, bao ngô nặng hơn bao gạo 9 kg. Bao ngô cân nặng bao nhiêu ki-lô-gam?
\shortans{55}
\loigiai{
Bao ngô nặng: $46 + 9 = 55$ kg.
}
\end{ex}

\begin{ex}
Tìm một số biết lấy số đó trừ đi 27 thì được 38.
\shortans{65}
\loigiai{
Số đó là: $38 + 27 = 65$.
}
\end{ex}

\begin{ex}
Tìm một số biết lấy số đó trừ đi 19 thì được 56.
\shortans{75}
\loigiai{
Số đó là: $56 + 19 = 75$.
}
\end{ex}

\begin{ex}
Tổng của số liền trước 48 và số liền sau 25 là bao nhiêu?
\shortans{73}
\loigiai{
Số liền trước 48 là 47, liền sau 25 là 26. Tổng: $47 + 26 = 73$.
}
\end{ex}

\begin{ex}
Một đoàn tàu chở hàng có 29 toa hàng khô và 18 toa hàng lạnh. Đoàn tàu có tất cả bao nhiêu toa hàng?
\shortans{47}
\loigiai{
Đoàn tàu có: $29 + 18 = 47$ toa.
}
\end{ex}

\begin{ex}
Tính: $19 + 29 + 39$.
\shortans{87}
\loigiai{
Ta có $19 + 29 = 48$, rồi $48 + 39 = 87$.
}
\end{ex}

\begin{ex}
Hai số có hiệu bằng 18, số trừ là 35. Số bị trừ là bao nhiêu?
\shortans{53}
\loigiai{
Số bị trừ: $18 + 35 = 53$.
}
\end{ex}

\begin{ex}
Tính tổng của số lớn nhất có hai chữ số khác nhau và số nhỏ nhất có một chữ số khác 0 (là 98 và 1).
\shortans{99}
\loigiai{
Tổng là: $98 + 1 = 99$.
}
\end{ex}
